\chapter{La traque}
\label{chap:la-traque}

La réussite de sa première épreuve avait changé certaines choses pour Khaer
Ghal. Il restait un enfant et personne au village ne l'aurait encore considéré
comme un guerrier, mais il avait désormais prouvé qu'il était capable de se
défendre seul. La longue canine qu'il portait à l'oreille gauche en témoignait
à chaque fois qu'il croisait un autre membre du clan, tout comme le contour
irrégulier du cartilage, définitivement marqué par la morsure qui la lui avait
fait gagner.

Peu à peu, les limites qu'on lui imposait s'éloignèrent. On lui permit
d'accompagner les chasseurs dans les environs du village, d'abord pour de
courtes sorties, puis pendant des journées entières. Quelque temps plus tard,
il fut également autorisé à suivre certains guerriers lors de leurs patrouilles.

Ces excursions constituaient une nouvelle étape de son apprentissage. Dans les
bois, Khaer Ghal devait reconnaître les traces des animaux et les distinguer de
celles laissées par des voyageurs, apprendre à progresser sans ralentir le
groupe et repérer les signes d'un passage récent. Les guerriers lui montraient
comment choisir un itinéraire permettant de rester dissimulé, comment estimer
le nombre d'individus ayant emprunté une piste ou encore comment reconnaître,
à la profondeur et à la régularité d'une empreinte, si celui qui l'avait
laissée marchait chargé.

Khaer Ghal aimait particulièrement ces patrouilles. Elles l'emmenaient toujours
un peu plus loin de la palissade et lui permettaient de découvrir un territoire
qu'il ne connaissait jusque-là qu'à travers les récits des chasseurs. Il
apprenait à reconnaître certains ruisseaux, des rochers aux formes singulières,
les passages empruntés par le gibier ou les arbres qui pouvaient servir de
repères lorsque le ciel disparaissait derrière les feuillages.

Il aimait surtout observer.

Trouver une trace presque effacée, comprendre ce qui l'avait produite puis
constater quelques centaines de pas plus loin qu'il avait vu juste lui
procurait une satisfaction différente de celle du combat. Il n'y avait ni
spectateurs ni cris d'encouragement, seulement un détail aperçu dans la boue,
une branche cassée ou quelques feuilles retournées. Pourtant, chaque fois qu'il
parvenait à remarquer quelque chose avant qu'un adulte ne le lui montre, il
avait l'impression d'avoir progressé.

À dix ans, il pouvait désormais passer une journée entière loin du village et
ne retrouver la palissade qu'à la tombée de la nuit, couvert de terre,
fatigué, parfois trempé, mais presque toujours impatient de repartir. Ce fut au
cours de l'une de ces sorties que son apprentissage prit une direction
différente.

La patrouille avait quitté le village tôt ce matin-là. Khaer Ghal accompagnait
plusieurs guerriers expérimentés et avançait, comme à son habitude, légèrement
derrière eux. Il n'était pas là pour choisir leur route ni pour participer aux
décisions ; son rôle consistait à regarder, écouter et retenir tout ce qu'il
pouvait.

Depuis plusieurs heures, rien n'avait troublé leur progression. Les Orcs
suivaient des passages qu'ils connaissaient bien, s'arrêtant parfois pour
examiner une piste ou contourner une zone plus difficile. Après une nuit
humide, le sol conservait particulièrement bien les empreintes.

Ce fut Khaer Ghal qui remarqua les premières.

À quelques pas du sentier qu'ils suivaient, une portion de terre meuble avait
été profondément marquée. Il s'arrêta, revint légèrement en arrière puis
s'accroupit. Plusieurs traces de sabots se superposaient et, entre elles, deux
sillons parallèles avaient creusé la boue.

Des roues.

Il suivit les marques du regard, puis examina les empreintes autour d'elles.
Certaines étaient déjà déformées par les passages suivants, mais elles étaient
trop nombreuses pour appartenir à quelques voyageurs isolés. Un guerrier
remarqua son arrêt et vint s'accroupir à côté de lui.

Khaer Ghal désigna les ornières.

— Plusieurs chariots.

L'Orc ne répondit pas immédiatement. Il observa le sol, écarta quelques
feuilles du bout des doigts puis suivit lui aussi la direction de la piste.

— Combien ?

Khaer Ghal regarda de nouveau les traces. Il aurait pu répondre qu'il n'en
savait rien. À la place, il essaya.

— Plus de trois.

Le guerrier examina encore le terrain avant de lui adresser un bref signe de
tête.

— Au moins quatre.

Ce n'était pas tout à fait juste, mais Khaer Ghal avait trouvé la piste. Les
autres furent rappelés et la patrouille abandonna aussitôt son itinéraire
initial.

Ils avaient trouvé une caravane.

\scenebreak

La piste n'était pas difficile à suivre. Les roues avaient profondément marqué
les portions humides de la route et les sabots des animaux laissaient derrière
eux une multitude d'empreintes. Pourtant, à mesure que les Orcs progressaient,
leur manière de se déplacer changea.

Les conversations cessèrent. Les guerriers quittèrent le chemin pour avancer
sous le couvert des arbres, suffisamment près pour ne pas perdre la piste mais
assez loin pour ne pas apparaître à découvert. Khaer Ghal les imita, attentif
à l'endroit où il posait les pieds. Une branche cassée sous une botte ou le
bruissement trop brutal d'un buisson suffisait désormais à attirer un regard
désapprobateur.

Au bout d'un moment, il entendit quelque chose : d'abord un grondement
irrégulier, presque confondu avec les bruits de la forêt, puis le grincement
caractéristique d'une roue et le claquement sec d'un harnais. Quelques voix
leur parvinrent ensuite, trop lointaines pour distinguer les paroles.

Les Orcs ralentirent encore et finirent par atteindre un endroit d'où ils
pouvaient observer la route entre les arbres. Khaer Ghal s'accroupit derrière
un tronc et découvrit pour la première fois le groupe qu'ils suivaient depuis
plusieurs heures.

Quatre chariots avançaient lentement sur le chemin. Des chevaux et quelques
bêtes de bât tiraient les véhicules chargés de caisses, de sacs et de ballots.
Des hommes et des femmes marchaient à leurs côtés tandis que d'autres
voyageaient assis sur les chariots. Plusieurs portaient des armes et deux
d'entre eux surveillaient régulièrement les bois.

Khaer Ghal les compta, puis recommença.

Il y avait beaucoup plus de monde que dans leur patrouille.

Il observa l'un des guerriers, certain qu'il allait donner le signal du retour.
À sa surprise, celui-ci désigna plutôt un Orc qui disparut aussitôt entre les
arbres dans la direction du village. Les autres restèrent et, pendant plusieurs
heures, suivirent la caravane à distance.

Khaer Ghal découvrit ce jour-là une autre forme de chasse. Il ne s'agissait
plus de suivre une bête jusqu'à son terrier ou d'attendre qu'un animal
s'approche suffisamment pour l'abattre. Cette fois, les guerriers observaient
des voyageurs. Ils comptaient ceux qui portaient des armes, surveillaient
l'allure des chariots, repéraient les endroits où le convoi ralentissait et
étudiaient le terrain autour de la route.

Khaer Ghal fit de même. Il ne comprenait pas encore tout ce qu'ils cherchaient,
mais il commençait à voir certaines choses : un chariot plus lourd que les
autres, un garde qui s'éloignait régulièrement du groupe, la manière dont les
voyageurs se resserraient lorsque le chemin devenait plus étroit. Il remarqua
aussi qu'en fin de journée les hommes armés observaient moins souvent les bois
et que les conducteurs semblaient surtout préoccupés par l'endroit où ils
pourraient s'arrêter pour la nuit.

Lorsque les renforts du village les rejoignirent, le soleil avait déjà commencé
à descendre entre les arbres. Les adultes se réunirent à l'écart ; Khaer Ghal
n'entendit que quelques mots et ne participa évidemment pas à leur décision.
Lorsqu'ils se dispersèrent, un guerrier lui ordonna simplement de rester en
retrait et de n'avancer que si on le lui demandait.

Khaer Ghal obéit.

\scenebreak

L'attaque commença peu avant la tombée de la nuit.

Les premiers Orcs surgirent des bois devant la caravane tandis qu'un second
groupe apparaissait derrière elle. Des cris éclatèrent presque immédiatement.
Un conducteur tira brutalement sur ses rênes ; derrière lui, un second chariot
tenta de l'éviter, quitta partiellement la route et s'immobilisa de travers
entre deux arbres.

Les voyageurs armés réagirent presque aussitôt. Certains se regroupèrent devant
ceux qui ne pouvaient pas combattre, tandis que d'autres tentèrent de repousser
les premiers assaillants. La confusion transforma rapidement la route en un
enchevêtrement de chevaux, de chariots et de silhouettes courant dans toutes
les directions.

Khaer Ghal resta sous les arbres. Il en avait reçu l'ordre.

Depuis sa position, il voyait les guerriers du village agir d'une manière qu'il
connaissait sans l'avoir jamais réellement observée. Ils avançaient par petits
groupes, profitaient des chariots pour isoler leurs adversaires et ne
s'attardaient jamais inutilement sur quelqu'un qui cessait de représenter une
menace.

Un garde leva son épée et tenta de tenir seul le passage entre deux véhicules.
Un Orc le força à reculer pendant qu'un second le contournait. Plus loin, un
voyageur jeta son arme au sol et leva les mains.

Personne ne le frappa.

Khaer Ghal suivit la scène des yeux. Quelques instants plus tard, un autre
homme fit de même et il comprit alors la règle appliquée par les guerriers :
ceux qui continuaient à se battre étaient combattus ; ceux qui se rendaient
étaient capturés.

La distinction lui parut naturelle. Elle ressemblait à beaucoup de choses qu'on
lui avait enseignées : une menace devait être éliminée ; quelqu'un qui ne
représentait plus un danger pouvait être utile autrement.

Lorsque les derniers voyageurs armés cessèrent de résister, un guerrier fit
signe à Khaer Ghal de le rejoindre. Plusieurs survivants avaient été regroupés
à l'écart de la route et débarrassés de leurs armes. On lui ordonna de rester
près d'eux avec un autre Orc pendant que les adultes inspectaient les chariots.

Il prit sa place et fit ce qu'on attendait de lui. Les prisonniers parlaient
peu ; certains regardaient les Orcs avec colère, d'autres fixaient simplement
le sol. Khaer Ghal les surveilla comme on le lui avait demandé, sans chercher
à leur adresser la parole. Pour lui, ils n'étaient pas fondamentalement
différents des prisonniers qu'il avait déjà vus au village.

La seule différence était qu'il savait désormais comment ils y arrivaient.

\scenebreak

Lorsque tout danger eut disparu, l'agitation changea progressivement de nature.
Les guerriers cessèrent de surveiller les bois et commencèrent à parcourir les
chariots. Des caisses furent ouvertes, des sacs déplacés et les marchandises
étalées sur la route afin de décider de ce qui méritait d'être rapporté.

Khaer Ghal fut bientôt relevé de sa surveillance et rejoignit les autres. Il
observa avec attention leur manière de procéder. Les objets précieux étaient
mis de côté, mais ils n'étaient pas les seuls à intéresser les Orcs : une bonne
hache pouvait avoir davantage de valeur pour le village qu'un bijou.

Les outils furent donc rassemblés avec les armes, les provisions encore
utilisables réparties dans des sacs et même certaines pièces de tissu
récupérées. Un guerrier inspecta longuement une roue endommagée avant d'en
retirer plusieurs éléments métalliques ; un autre vida une caisse pour
conserver uniquement les récipients qu'elle contenait.

Khaer Ghal reconnut alors certains objets. Pas ceux qui se trouvaient devant
lui, mais leur équivalent : les couvertures utilisées dans certaines maisons,
les outils de l'atelier, les sacs dans lesquels on conservait le grain, une
boucle de harnais semblable à celles qu'il avait vues près des enclos. Il avait
toujours su que les guerriers rapportaient des choses de leurs patrouilles. Il
n'avait simplement jamais réfléchi à l'endroit précis où ils les avaient
prises.

Un peu plus loin, l'un des chariots trop endommagés pour être ramené fut vidé
de tout ce qui pouvait encore servir. Lorsque les Orcs eurent terminé, ils y
mirent le feu. Les flammes gagnèrent lentement le bois sec et une colonne de
fumée monta entre les arbres, traversée par les derniers rayons du soleil.

Les survivants capables de marcher furent rassemblés pour le retour. Ils
seraient conduits au village, où certains travailleraient pour le clan tandis
que d'autres pourraient un jour être envoyés dans l'arène. Les quelques
blessés dont l'état ne permettait pas le voyage ne repartiraient pas avec eux.

Khaer Ghal ne posa aucune question.

Il avait grandi parmi ces Orcs. Depuis son enfance, il voyait arriver au village
des armes, des outils, des marchandises et parfois des prisonniers. Il savait
que les guerriers attaquaient ceux qui s'aventuraient sur leur territoire ou
qui transportaient quelque chose dont le clan pouvait avoir besoin. Rien de ce
qu'il venait de voir ne contredisait véritablement ce qu'on lui avait appris.

Ce n'était donc pas encore du doute.

Seulement une différence.